\section{Controlled rare-subgroup prediction pilot}
\label{sec:pilot}

We compare four fixed scalings on six feasible two-binary-feature designs. Each
design has $n = 400$ observations, common prevalence $p = 1/2$, and rare
prevalence $q \in \{0.05,0.01\}$. The joint law is independent, nested ($J = 1$
implies $W = 1$), or negatively associated at the disjoint boundary
$\Prob(W = 1,J = 1) = 0$. Training cell counts equal their population
expectations, giving 20 or 4 rare observations. The signals are $(1,0)$,
$(0,1)$, $(1,1)$, and $(1,-1)$, with iid Gaussian observation noise of variance
one. Population centering and scales are known, the centered-model intercept is
fixed at zero, and coefficients start at zero. Repeated visits to a training
observation reuse its noise. This experiment concerns fixed-data sampling,
separate from the fresh-sample joint limit in Section~\ref{sec:budget}.

For coordinate variance $v_j$, the feature divisors are $1$, $\sqrt{v_j}$,
$v_j$, and $\sqrt{v_j + 0.01}$. We use full-batch GD and replacement SGD with
$B \in \{1,20\}$. Budgets of $M \in \{400,2000,8000\}$ processed observations
give $M/400$ GD updates and $M/B$ SGD updates. Every design, signal, method,
scaling, and budget receives separate constant-step tuning over the same nine
factors $a \in \{0.001,0.003,0.01,0.03,0.1,0.3,0.7,1.5,2.5\}$, with
$\eta = a/\lambda_{\max}(S^{-1}CS^{-1})$. The objective is mean rare-subgroup
risk across 32 tuning noise realizations. Evaluation uses 256 disjoint noise
realizations, paired across methods and scalings. This is oracle tuning with the
true signal and test law, rather than validation-based tuning. Nominal tuning
effort is $9 \cdot 32 \cdot M$ processed observations per comparison.

For original-scale error $e = \beta - \beta^*$, let $V_D = \E[ee^\top\mid D]$,
where $D$ is the fixed noisy training dataset and the expectation integrates
replacement-index sampling. We report excess squared prediction errors
\[
 R_D = \operatorname{tr}(CV_D),\qquad R_{D,R}
     = \operatorname{tr}(C_RV_D),\qquad C_R
     = \E[xx^\top\mid J=1].
\]
These reported risks omit the factor $1/2$ in the training loss and exclude
irreducible test noise. The conditional second moment $C_R$ includes the nonzero
conditional feature mean. An isolated rare-coefficient error of one therefore
contributes $q(1 - q)$ to aggregate risk and $(1 - q)^2$ to rare risk, a ratio
of 99 at $q = 0.01$.

Conditional means and second moments are computed exactly by powering the affine
replacement operator of Section~\ref{sec:sampling}, augmented by a constant
coordinate. This retains the empirical score and curvature--noise cross terms.
Independent enumeration of ordered batches and two-update paths, direct
prediction-risk calculations, and a separate moment recurrence check agree to
numerical precision. Monte Carlo standard errors (MCSEs) below concern training
noise only, conditional on the fixed designs and selected steps. They exclude
variation in tuning, random designs, and estimated preprocessing.

\begin{table}[t]
 \centering
 \small
 \setlength{\tabcolsep}{4pt}
 \begin{tabular}{lrllll}
  \toprule
  Method        &  $M$ & None          & Std. dev.     & Variance      & Damped        \\
  \midrule
  GD            &  400 & .2954 (.0013) & .1565 (.0026) & .1948 (.0024) & .1704 (.0026) \\
  GD            & 2000 & .2028 (.0023) & .1140 (.0031) & .1415 (.0037) & .1161 (.0035) \\
  GD            & 8000 & .1022 (.0022) & .1152 (.0034) & .1308 (.0046) & .1190 (.0033) \\
  SGD, $B = 1$  &  400 & .2275 (.0024) & .2119 (.0027) & .2554 (.0028) & .2228 (.0018) \\
  SGD, $B = 1$  & 2000 & .1384 (.0035) & .1565 (.0045) & .1914 (.0045) & .1506 (.0033) \\
  SGD, $B = 1$  & 8000 & .1257 (.0036) & .1795 (.0060) & .1995 (.0059) & .1662 (.0058) \\
  SGD, $B = 20$ &  400 & .2145 (.0017) & .2213 (.0023) & .2594 (.0022) & .1916 (.0024) \\
  SGD, $B = 20$ & 2000 & .1265 (.0022) & .1513 (.0035) & .1798 (.0030) & .1282 (.0029) \\
  SGD, $B = 20$ & 8000 & .1058 (.0030) & .1290 (.0029) & .1522 (.0033) & .1180 (.0037) \\
  \bottomrule
 \end{tabular}
 \caption{
  Rare-presence excess squared prediction error, with MCSE in parentheses. Each
  entry averages equally over all six designs and four signals. MCSEs preserve
  the pairing across signals within a design. Equal weighting is a descriptive
  choice for this finite grid.
 }
 \label{tab:pilot}
\end{table}

Table~\ref{tab:pilot} shows substantial differences and changes in their
direction. At $q = 0.01$ with independent features and $M = 400$,
standard-deviation scaling changes rare risk relative to no scaling by
$-0.58017$ (paired MCSE $0.01517$) for the rare-only signal and $+0.14895$
($0.01301$) for the common-only signal. This budget permits only one GD update.
Dependence also reverses equal-sign and opposite-sign comparisons. Across all 24
cases, the average GD advantage of standard-deviation scaling at $M = 400$
reverses by $M = 8000$. The corresponding reversal for $B = 1$ SGD occurs by
$M = 2000$. These are examples of finite-budget preconditioning and
noise-fitting tradeoffs, without a general ordering of the scaling rules.

The finite tuning grid is a material limitation. Its smallest factor is selected
in 195 of 864 evaluated settings and its largest in nine. All selected endpoints
are finite, but nine GD selections are asymptotically unstable and are evaluated
after only one or five updates. Of 7776 candidate endpoints, 2160 have
homogeneous second-moment spectral radius at least one and 1028 have nonfinite
or invalid moments; these categories overlap. All outcomes are retained. The
grid was fixed before evaluation and was not expanded after observing these
boundary selections. The repository's \texttt{notes/rare-subgroup-pilot.md}
records the full protocol, aggregate risks, casewise contrasts, limitations, and
reproduction commands. The generated JSON retains coefficients, covariance,
risks, counts, scales, selected rates, seeds, and software provenance. The
experiment supplies a controlled illustration, without establishing robustness
to a wider tuning search, estimated preprocessing, or real-data applications.
